% ============================================================================
% PATRON ENTRY: THE RIVEN PATH (PASSION & POWER)
% ============================================================================

\subsection{The Riven Path (Passion \& Power)}

\subsubsection*{Lore}
The Riven Path does not command. It \textit{offers}. It is the heat in the blood, the fire in the belly, the voice that whispers \textit{``Take what is yours.''} Those who walk this path learn that power is not given --- it is \textit{seized}. Chains are not broken by asking nicely. They are shattered by will.

Followers of the Riven Path are conquerors, liberators, and survivors. They serve no master --- not even the Path itself. When a tyrant rises, they do not negotiate --- they \textit{overthrow}. When a wall stands in their way, they do not climb --- they \textit{break}. When a rival challenges them, they do not yield --- they \textit{prove}.

Among the peoples of Fate's Edge, the Riven Path is known by names that frighten children and inspire revolutionaries. In Silkstrand, it is the \textit{Cut Purse's Blessing}, whispered by thieves who take from those who stole first. In Vhasia, it is the \textit{Gorget-Breaker}, a philosophy of the duelist who refuses to wear armor. In Zakov, it is the \textit{Tide-That-Eats-The-Shore}, a slow, patient consumption that reshapes coastlines. Wherever the oppressed dream of rising, the Riven Path waits.

\begin{quote}
``The chain that binds you is the chain you refuse to break. The hand that holds you down is the hand you have not yet bitten.''
\end{quote}

\subsubsection*{Domain Focus}
\begin{itemize}
  \item \textbf{Liberation Through Power:} freedom as the right of the strong, justice as the will of the victorious
  \item \textbf{Passion as Fuel:} anger, desire, fear, and love --- all are fires to be stoked and aimed
  \item \textbf{Crushing Dominance:} victory through overwhelming force, intimidation, and the willing display of power
  \item \textbf{Self-Determination:} no chains, no masters, no debt that is not paid in blood
\end{itemize}

\subsubsection*{Thiasos and Codex (Runekeeper Options)}

A Runekeeper sworn to the Riven Path does not bind a living creature of loyalty or service --- they bind a fragment of broken power. Their \textbf{Thiasos} is a \textit{cracked crystal}: a focus that was once whole, now broken, now \textit{more}. It burns warm when they are angry and grows hot when they are in danger. It requires no food, but they must feed it a drop of their own blood each day (mark 1 Fatigue once per week if neglected). If neglected for a week, it becomes \textit{Compromised} (it goes cold, imposing --1 die on all Intimidation and Command rolls until they feed it).

The \textbf{Codex} of the Riven Path is never a book of prayers. It is a \textit{manifesto} --- a collection of maxims, battle cries, and personal victories. It cannot be read in the usual way; the follower must \textit{earn} each Rite by proving their worth through combat, theft, or domination. Because the Path's truths are written in acts of dominance, the Codex requires upkeep: the Runekeeper must perform one act of dominance per week (win a spar, intimidate someone, break something that held them back). If neglected, it becomes \textit{Compromised} (the Rites become illegible until you win a meaningful victory).

\subsubsection*{Patron's Gift (Imbuement)}

Once per scene as an action, a Runekeeper of the Riven Path may touch a held weapon or focus and whisper a command of breaking. For the rest of the scene, the item grants \textbf{+1 die} to \textit{Melee} or \textit{Intimidation} rolls, and the first time each scene the bearer would be disarmed or restrained, they may reroll the resistance check (the Path's fury surges). After the scene, the user marks \textbf{+1 Obligation} to the Riven Path.

\subsubsection*{Rites of the Riven Path}

\paragraph{Rite of Feasting on Fear} (Basic, 4 XP) \hfill \\
\emph{Instant; Self; No.} \\
\textbf{Tags:} \texttt{[DRAIN]}, \texttt{[FEAR]} \\
\textbf{Materials:} A target that is actively frightened or in your power. \\
\textbf{Effect:} When you cause a creature to become \textit{Frightened} or gain a fear-based Condition, you immediately clear 1 Fatigue. The terror of others sustains you. \\
\textbf{Push It:} Clear 2 Fatigue instead of 1. Mark 1 SB (Hearts) as the fear becomes palpable to nearby observers. Cost: Mark +1 Obligation (total +2) or Mark +1 Fatigue. \\
\textbf{Cost:} Mark +1 Obligation. \\
\textit{Requires:} Familiar (\textit{Invoke:} 1 Boon). \\
\textbf{Invoke:} 1 action. \\
\textbf{Duration:} Instant.

\paragraph{Rite of the Crushing Grip} (Basic, 5 XP) \hfill \\
\emph{Instant; Near; Yes.} \\
\textbf{Tags:} \texttt{[FORCE]}, \texttt{[STRANGLE]} \\
\textbf{Materials:} A closed hand, reaching toward the target, fingers slowly tightening. \\
\textbf{Effect:} One target within Near must test \textit{Body + Athletics} vs. your \textit{Spirit + Lore}. On a failure, they take 1 Harm and cannot speak, breathe, or cast spells with verbal components until they break free (test again as an action each round). You must maintain concentration (cannot take other actions) or the effect ends. \\
\textbf{Push It:} The target also suffers --1 die to all actions while gripped. Cost: Mark +1 Obligation (total +2) or Mark +1 Fatigue. \\
\textbf{Cost:} Mark +1 Obligation. \\
\textit{Requires:} Familiar (\textit{Invoke:} 1 Boon). \\
\textbf{Invoke:} 1 action. \\
\textbf{Duration:} Scene.

\paragraph{Rite of the Cracking Lightning} (Advanced, 9 XP) \hfill \\
\emph{Instant; Close; Yes.} \\
\textbf{Tags:} \texttt{[STORM]}, \texttt{[STRIKE]}, \texttt{[AREA]} \\
\textbf{Materials:} A raised hand, fingers spread, as if holding a sphere of static. \\
\textbf{Effect:} A cone of arcing lightning lashes from your hand. All targets in Close range must test \textit{Body + Athletics} (DV 4) or take \textbf{Harm 2}. Those who succeed take Harm 1 instead. \\
\textbf{Push It:} The range extends to Near. Cost: Mark +1 Obligation (total +2) or Mark +1 Fatigue. \\
\textbf{Cost:} Mark +1 Obligation. \\
\textit{Requires:} Familiar + Codex (\textit{Invoke:} 1 Boon). \\
\textbf{Invoke:} 1 action. \\
\textbf{Duration:} Instant.

\paragraph{Rite of the Battle Rage} (Advanced, 8 XP) \hfill \\
\emph{Scene; Self; No.} \\
\textbf{Tags:} \texttt{[PASSION]}, \texttt{[BOOST]} \\
\textbf{Materials:} A moment of silent, building fury. \\
\textbf{Effect:} For the next three exchanges, you gain \textbf{+2 dice} to all melee attacks, ignore the first Harm 1 each round, and ignore all Fatigue penalties. However, your Position becomes \textit{Desperate} (re-roll one success) for the duration --- you are reckless, not careful. When the rage ends, mark 1 Fatigue. \\
\textbf{Push It:} The rage lasts for the full scene. Mark 2 Obligation and mark 1 SB (Diamonds) as collateral damage. Cost: Mark +1 Obligation (total +3) or Mark +1 Fatigue. \\
\textbf{Cost:} Mark +1 Obligation. \\
\textit{Requires:} Familiar + Codex (\textit{Invoke:} 1 Boon). \\
\textbf{Invoke:} 1 action. \\
\textbf{Duration:} 3 exchanges.

\paragraph{Rite of the Crushing Dark} (Master, 13 XP) \hfill \\
\emph{Instant; Zone; Yes.} \\
\textbf{Tags:} \texttt{[FORCE]}, \texttt{[AREA]}, \texttt{[FEAR]}, \texttt{[PRESSURE]} \\
\textbf{Materials:} A clenched fist and a whispered promise of pain. \\
\textbf{Effect:} Invisible pressure descends on all enemies within the zone. Each target must test \textit{Body + Endurance} (DV 4) or suffer \textbf{Harm 1} (crushing) and be \textbf{knocked prone}. Even on a successful save, the target suffers \textbf{1 Fatigue} from the sheer pressure and cannot voluntarily move closer to you until the end of their next turn (they are held at bay). \\
\textbf{Push It:} The pressure intensifies to a bone-crushing degree. All enemies in the zone who fail their save instead suffer \textbf{Harm 2} and are \textbf{Stunned} (lose their next action). Those who succeed still suffer \textbf{1 Fatigue} and are pushed back one range band. Cost: Mark +1 Obligation (total +3) or Mark +1 Fatigue. \\
\textbf{Cost:} Mark +2 Obligation. \\
\textit{Requires:} Familiar + Codex + Tier III (\textit{Invoke:} 2 Boons). \\
\textbf{Invoke:} Extended action. \\
\textbf{Duration:} Instant.

\paragraph{Rite of the Riven Fury} (Master, 13 XP) \hfill \\
\emph{Scene; Self; No.} \\
\textbf{Tags:} \texttt{[STRIKE]}, \texttt{[BRUTAL]} \\
\textbf{Materials:} A blade or bo staff held in a defensive stance. \\
\textbf{Effect:} For the rest of the scene, you gain +1 defense dice. Once per round, when you are targeted by a melee attack, you may make a counter-attack against the attacker. You reduce the Harm you would have taken from the triggering attack by 1 (minimum 0). The counter-attack gains +1 die and deals +1 Harm on a hit. \\
\textbf{Push It:} The counter-attack also ignores the attacker's Armor. Cost: Mark +1 Obligation (total +3) or Mark +1 Fatigue. \\
\textbf{Cost:} Mark +2 Obligation. \\
\textit{Requires:} Familiar + Codex (\textit{Invoke:} 1 Boon). \\
\textbf{Invoke:} 1 action. \\
\textbf{Duration:} Scene.

\subsubsection*{Cantors \& Cults of the Riven Path}

\textbf{The Cantor's Song.} A Cantor of the Riven Path sings in \textit{war cries} --- rhythmic chants that build into ecstatic frenzy, that drive soldiers to charge and cowards to flee. Their voice is used to terrify enemies, to rally allies on the brink of breaking, or simply to announce their presence in a room. A Cantor may learn any Low Rite of the Riven Path as a Song (\textit{Performance + Lore}, resolves next turn). Corruption accumulates when they use their voice for restraint, for mercy, or to calm a situation they could have escalated.

\textbf{The Cult of the Chain-Breakers.} The Riven Path's cults gather in places where chains have been broken --- former slave quarters, prisons after a riot, the ruins of a tyrant's palace. The Chain-Breakers are a fellowship of escaped captives, freed debtors, and those who have chosen to owe nothing to anyone. A ritual of the Cult involves each member speaking aloud a chain they still wear --- a fear, a debt, a relationship that controls them --- and then breaking a physical chain (rope, iron links, a leather cord) in front of witnesses. The broken pieces are thrown into a fire. Their creed is shouted, never whispered: \textit{``I owe nothing. I serve no one. I am my own.''} Outsiders are not welcome unless they bring a chain to break. Those who break a chain with sincerity may stay for the feast. Those who hesitate are watched until they leave.

\subsubsection*{Witchcraft of the Riven Path (Lay / Unmaking)}

A witch who walks with the Riven Path is a \textit{chain-shatterer} or \textit{oath-breaker} --- one who crafts through the art of destruction. Their magic works through the materials of power seized: a hammer that breaks locks, a key that opens any door once (then crumbles), or a brand that leaves a mark that cannot be healed. In Zakov, a pirate captain might use a \textit{mutineer's compass} that always points to the nearest source of fear, but a witch of the Riven Path would instead carry the fingerbone of a captain they killed, and the bone points to that captain's hidden treasure.

\textbf{The Witch's Tool: The Chain-Breaker's Hammer.} A blacksmith's hammer with a cracked handle, wrapped in leather from a slaver's belt. Once per session, the witch may strike a lock, a chain, or a sealed document with the hammer. The target does not break --- it \textit{ceases to be locked}. The mechanism is not damaged; it simply no longer functions as a barrier. The lock remains, but it will not close. The chain remains, but it will not hold. The seal remains, but it will not seal. The effect is permanent. The hammer may be used on magical seals, but the witch must succeed on a \textit{Spirit + Lore} test (DV 4) or the hammer cracks further (treat as \textit{Compromised} until repaired with a drop of blood).

\textbf{A Hedge Gift: The Unbound Step (Basic, 4 XP).} Once per scene, the witch may ignore one physical restraint --- ropes, manacles, a grappling hand, a closed door --- as if it were not there for a single exchange. They do not break it; they simply \textit{are not bound by it}. The rope still exists. The manacles are still closed. But the witch is standing three feet away. The GM gains 1 SB (the world takes a moment to notice the absence). This gift cannot be used again until the gained SB is spent or the scene ends.

\subsubsection*{Corruption of the Riven Path}

\begin{tblr}{
  colspec = {Q[c,1cm] X[2.5] X[2.5]},
  rowhead = 1,
  row{odd} = {bg=gray!10},
  row{1} = {bg=gray!30, font=\bfseries},
  hlines,
}
Tier & Benefit & Cost / Quirk \\
1 & \textbf{Predator's Focus:} +1 die to Intimidation and Command rolls. & \textbf{Hunger:} --1 die to resist acting on strong emotion (anger, lust, greed, vengeance). \\
2 & \textbf{Blood Scent:} Once/scene, automatically locate the nearest wounded creature within Far range. & \textbf{Mercy Blind:} --1 die to show compassion, restraint, or forgiveness. \\
3 & \textbf{Unstoppable:} +2 dice to resist being grappled, bound, or restrained. & \textbf{Reckless:} --1 die to defensive actions (you favor attack over parry). \\
4 & \textbf{Fear Eater:} Once/scene, when you cause a creature to become Frightened, gain 1 Boon. & \textbf{Isolation:} --1 die to social rolls with anyone you have not intimidated. \\
5 & \textbf{Crack the Chain:} Once/session, automatically break any mundane restraint (rope, manacles, cage, locked door). & \textbf{Hunted:} A rival from your past appears. The GM gains 1 SB per session to spend on their pursuit. \\
6+ & \textbf{Unbound:} Once/session, ignore all penalties from Fatigue and Harm for one scene. & \textbf{The Rift Claims You:} Mark +3 Obligation. You may no longer take any action that requires submission, surrender, or yielding. You cannot be reasoned with --- only fought or fled. \\
\end{tblr}

\subsection*{Playstyle Notes}
The Riven Path is the ideal patron for players who want to play an aggressive, self-determined character --- a warlord, a revolutionary, a dark knight, or simply someone who refuses to bow. Its Rites reward aggression, intimidation, and the willingness to take what you want. The cost is isolation --- the power that lets you break chains also breaks relationships, trust, and your ability to see others as anything but obstacles or tools. Intrusions often take the form of a rival who has also broken their chains and now wants yours, a debt you cannot pay (because you broke the contract), or a situation where mercy would be easier than violence --- forcing you to choose between your principles and your mission. Ideal for sith analogues, warlocks, berserkers, and any character who believes that freedom is the right of the strong and that weakness is a choice.

\noindent\textbf{Emphasizes}
\begin{itemize}
  \item \textbf{Liberation Through Power:} freedom as the right of the strong
  \item \textbf{Passion as Fuel:} anger, desire, fear --- all are fires to be stoked
  \item \textbf{Crushing Dominance:} victory through overwhelming force
  \item \textbf{Self-Determination:} no chains, no masters, no debt without blood
  \item \textbf{The Breaking Wave:} the joy of watching something that held you crumble
\end{itemize}

% --- Monastic Tradition: Riven-Fist (The Riven Path) ---
\begin{tcolorbox}[colback=black!5,colframe=red!80!black,title=\textbf{Riven-Fist (The Riven Path)},breakable]
  \emph{The chain that binds you is the chain you refuse to break. The hand that holds you down is the hand you have not yet bitten.}

  \vspace{2mm}
  \noindent\textbf{Prerequisites:} Body 3+, Spirit 2+, Rite of Feasting on Fear or Cantor of the Riven Path.

  \vspace{2mm}
  \noindent\textbf{Debt-Resistant Frame.} The Riven Path does not forgive debts --- it \textit{collects}. Monks of the Riven-Fist gain \textbf{+1 die} to resist any magical debt, geas, or contract that would compel them to submit or yield. If they succeed, the creditor feels a momentary chill, as if they have offended something hungry --- and the GM gains 1 SB (the Path notes a new target).

  \vspace{2mm}
  \noindent\textbf{Techniques (purchased as Talents):}

  \begin{tblr}{
    colspec = {X[0.7,l] X[1.8,l] X[1.1,l] X[1.4,l]},
    rowhead = 1,
    row{odd}={bg=gray!10},
    row{1}={bg=gray!30, font=\bfseries},
    hlines,
  }
  Technique & Effect & Cost & Req. \\
  \textbf{Feast on Fear} (Basic, 5 XP) & When you cause a creature to become Frightened, clear 1 Fatigue. The terror sustains you. & 1 Boon & Spirit 2 \\
  \textbf{Crushing Grip} (Basic, 5 XP) & Target within Near tests Body+Athletics vs. Spirit+Lore. On failure, 1 Harm and cannot speak or breathe until they break free. & 1 Fatigue & Body 3 \\
  \textbf{Cracking Lightning} (Advanced, 9 XP) & Close range cone. All targets test Body+Athletics (DV 4) or take Harm 2. Success = Harm 1. & 1 Boon & Spirit 3, Tier II \\
  \textbf{Battle Rage} (Advanced, 8 XP) & For 3 exchanges: +2 dice melee, ignore first Harm 1/round, ignore Fatigue penalties, Position Desperate. End: 1 Fatigue. & 1 Boon & Body 3, Tier II \\
  \textbf{Sundering Blade} (Advanced, 8 XP) & Next melee attack with light blade: +2 dice, ignore 1 Armor, +1 Harm. Crit = permanent Scar. & 1 Fatigue + 1 Boon & Tier II \\
  \end{tblr}

  \vspace{2mm}
  \noindent\textbf{Master Technique (Epic Talent, 12 XP):}
  \textit{Unbound.} Once per session, ignore all penalties from Fatigue and Harm for one scene. You feel no pain, no doubt, no hesitation --- only the ecstasy of action. After the scene, you gain a permanent \textbf{Riven Scar} (your eyes are slightly wrong --- the pupils are vertical slits for one heartbeat when you are angry). You may ignore one Obligation call from the Riven Path --- but the Path will remember, and the next call will be louder.
\end{tcolorbox}